%======================================================================%
\section{The Physical Universe}
\label{sec:physics}
%======================================================================%

Mathematical data of $\mathfrak{U}^\star$ becomes physical through mirror
quotient observability (Section~\ref{sec:observer}), observer collapse at
$\Mirror^7(\Void)$, and soldering on the Jordan--octonionic block.

%----------------------------------------------------------------------%
\subsection{Physical sector}
\label{sec:physics:sector}
%----------------------------------------------------------------------%

Physical configurations are observer-collapsed states on $\hat Q$
(Definitions~\ref{def:Qhat},~\ref{def:observer},~\ref{def:physical}).

\begin{definition}[Physical configurations]\label{def:physical}
A configuration is \emph{observable} if it lies in the $\sigma$-anti-invariant
sector of the mirror quotient $\hat Q=Q/\sigma$
(Definition~\ref{def:Qhat}). Fermions are cohomology classes in
$H^1(\hat Q,V_{16})^{\sigma=-1}$; each observer $|\omega_\iota\rangle$ is one
such class (Theorem~\ref{thm:observer-fp}). Observable states are those
surviving observer collapse $\Pi_{\sigma=-1}$
(Definition~\ref{def:heterotic-collapse}); $\sigma$-invariant mirror doubles
are quotient-non-observable.
\end{definition}

\begin{proposition}[No-go: P2 on the unquotiented coset]\label{prop:p2-nogo}
Void, mirror, and capacity balancing fix the \emph{number} of horizontal
family cells ($k=3$, Theorem~\ref{thm:partition}) and organize them by triality
on $\mathfrak{d}$, but they do \emph{not} derive chiral $\Spin(10)$ matter on
the unquotiented coset. The EIII matter directions
$\mathfrak{m}=\mathbf{16}_{-3}\oplus\overline{\mathbf{16}}_{+3}$
(Theorem~\ref{thm:universe}) are a mirror-symmetric pair
(Proposition~\ref{prop:mirror:involution}); selecting one chiral
$\mathbf{16}$ per family with no $\overline{\mathbf{16}}$ partners is
\emph{Postulate~P2}, not a consequence of stabilisation alone. A bare Yukawa
$\bar\psi\,\orderparam\,\psi$ with
$\orderparam\in\{\mathbf{16},\overline{\mathbf{16}}\}$ is not $H$-invariant on
the unquotiented coset.
\end{proposition}

\begin{remark}[Observer collapse eliminates P2]
Proposition~\ref{prop:p2-nogo} is not overridden: it records an honest limit
on \emph{postulated} chiral insertion atop the mirror-symmetric coset. The
elimination route \emph{redefines} the matter sector. Observer collapse
$\Pi_{\sigma=-1}$ together with Definition~\ref{def:physical} declares
physical fermions to be $\sigma$-anti-invariant classes on $\hat Q$; mirror
partners live in $H^1(\hat Q,V_{16})^{\sigma=+1}$ and are not observable.
Chirality is then a theorem of quotient observability, not an added input.
\end{remark}

\begin{theorem}[Chirality from observer collapse]\label{thm:chirality-derived}
Under Definition~\ref{def:physical} and Theorem~\ref{thm:observer-fp},
\begin{equation}\label{eq:chirality-derived}
  \dim H^1(\hat Q,V_{16})^{\sigma=-1}=3,
\end{equation}
with one chiral $\mathbf{16}$ of $\Spin(10)$ per observer index
$\iota\in\mathcal{I}$. No $\overline{\mathbf{16}}$ partners appear in the
physical sector because mirror partners are $\sigma$-invariant, not
anti-invariant. Hence $n_F=3$. This \textbf{replaces Postulate~P2}: chirality
is not postulated; it is the quotient definition of observability together
with the observer fixed-point count.
\end{theorem}

\begin{proof}
Theorem~\ref{thm:observer-fp} gives the anti-invariant dimension and the
bijection $\iota\mapsto|\omega_\iota\rangle$. Physical configurations are
anti-invariant by Definition~\ref{def:physical}; $\sigma$-invariant modes are
mirror doubles and quotient-non-observable. Cellular details in
Appendix~\ref{sec:appendix:quotient}.
\end{proof}

\begin{corollary}[Three families]\label{thm:three-families}
Theorem~\ref{thm:chirality-derived} yields $n_F=3$ chiral families without
importing heterotic axioms as primitive data~\cite{Witten1985}.
\end{corollary}

%----------------------------------------------------------------------%
\subsection{Spacetime from the climb}
\label{sec:physics:spacetime}
%----------------------------------------------------------------------%

Four-dimensional Lorentzian geometry is not a consequence of the $\E{6}$
coset metric alone. The mirror climb supplies a quaternionic clock block;
observer lifetime charts (Section~\ref{sec:observer:lifetime}) admit a
soldering map that imports signature into the contact layer.

\begin{definition}[Jordan--octonionic clock block]\label{def:spacetime:H2O}
Identify off-diagonal Jordan slots of $\jalg$ with the $2\times2$
Hermitian--octonionic model $H_2(\Oct)$. The determinant $\det$ carries
signature $(1,9)$ on $H_2(\Oct)\cong\R^{10}$. Restrict to the quaternionic
subalgebra $\Quat\subset\Oct$ fixed by the mirror involution at the
$4\to5$ Jordan exit (Theorem~\ref{thm:rung-table}). The embedded block
$H_2(\Quat)\subset H_2(\Oct)$ is four-dimensional; $\det|_{H_2(\Quat)}$
has Lorentzian signature $(-,+,+,+)$.
\end{definition}

\begin{definition}[Soldering]\label{def:soldering}
Let $\sigma_{\mathrm{sol}}:H_2(\Quat)\hookrightarrow\jalg$ be the canonical
embedding of Definition~\ref{def:spacetime:H2O}. The \emph{soldering map}
identifies four clock components $\orderparam^0,\ldots,\orderparam^3$ with
$\mathrm{Im}(\sigma_{\mathrm{sol}})$; the world manifold has $d=4$.
\end{definition}

\begin{lemma}[$F_4$ equivariance fixes the quaternionic embedding]
\label{lem:quat-F4}
Let $\iota_{\Quat}:\Quat\hookrightarrow\Oct$ be an associative subalgebra
embedding. Extend it to the Jordan--octonionic block
$H_2(\Quat)\hookrightarrow H_2(\Oct)\hookrightarrow\jalg$ by placing
quaternion entries in a single off-diagonal slot. Among such embeddings
compatible with
\begin{enumerate}[label=(\roman*),leftmargin=*]
  \item mirror functoriality on the climb, in particular the
    $4\to5$ Jordan exit that carries $\Mirror^3(\Void)=\Quat$ into
    $\Mirror^5(\Void)\cong\jalg$ (Theorem~\ref{thm:rung-table});
  \item $F_4=\mathrm{Aut}(\jalg)$ equivariance of the soldering map
    $\sigma_{\mathrm{sol}}$ (Definition~\ref{def:soldering}); and
  \item the above--below correspondence on the diagonal Cartan
    $\mathfrak{d}=\mathrm{span}\{e_1,e_2,e_3\}$
    (Definition~\ref{def:above-below}),
\end{enumerate}
the orbit under $F_4$ is a single $G_2$-conjugacy class, with
$G_2=\mathrm{Aut}(\Oct)\subset F_4$ the stabilizer of a generic
diagonal frame. Fixing an observer index $\iota\in\mathcal{I}$
(Definition~\ref{def:observer-index}) selects a unique representative
$\Quat_\iota\subset\Oct$ and hence a unique soldered block
$H_2(\Quat)\subset H_2(\Oct)$.
\end{lemma}

\begin{proof}[Proof sketch]
Associative subalgebras $\Quat\subset\Oct$ form one $G_2$-orbit
\cite{SpringerVeldkamp2000}. The Jordan envelope $\jalg$ carries three
off-diagonal octonion slots; $F_4$ permutes them while preserving
$\mathfrak{d}$. Mirror descent at $n=4\to5$ singles out the associative
subalgebra from rung $3$, so the $F_4$-admissible embeddings are those
$G_2$-conjugates compatible with capacity weights on $\mathfrak{d}$
(Proposition~\ref{prop:above-below}). Observer index $\iota$ labels one
diagonal direction $e_\iota$, breaking the residual triality and pinning
$\Quat_\iota$ uniquely.
\end{proof}

\begin{proposition}[Exceptional rungs preserve $H_2(\Quat)$ soldering]
\label{prop:soldering-transport}
On stabilized exceptional rungs $n=6,7,8$, the mirrors of
Proposition~\ref{prop:climb:exceptional} transport the soldered block
$H_2(\Quat)\subset H_2(\Oct)\hookrightarrow\jalg$ of
Definition~\ref{def:spacetime:H2O} without introducing a second clock sector:
\begin{enumerate}[label=(\roman*),leftmargin=*]
  \item $\E{8}$ ($n=6$): automorphism mirror couples the two $\Oct$ sectors of
    \eqref{eq:matrix}; the associative subalgebra $\Quat\subset\Oct$ from
    $\Mirror^3(\Void)$ is stabilized by $G_2=\mathrm{Aut}(\Oct)\subset\E{8}$,
    and $\det|_{H_2(\Quat)}$ retains Lorentzian signature $(-,+,+,+)$.
  \item $\E{8}\times\E{8}$ ($n=7$): product mirror $\sigma$ stabilizes the
    diagonal $\delta(\E{8})$ and the $\jalg$ chart; the soldering embedding
    $\sigma_{\mathrm{sol}}$ of Definition~\ref{def:soldering} is
    $\sigma$-equivariant on $H_2(\Quat)$, hence survives passage to
    $\hat Q=Q/\sigma$.
  \item $\E{6}$ ($n=8$): Cartan mirror $\Pi_{\mathrm{het}}$ collapses gauge data
    to $H=\Spin(10)\times U(1)$ while preserving $F_4=\mathrm{Aut}(\jalg)$
    equivariance (Lemma~\ref{lem:quat-F4}); the four clock components
    $\orderparam^0,\ldots,\orderparam^3$ remain the $\mathfrak{m}_{+1}$
    Goldstone block of Lemma~\ref{lem:goldstone-sigma}.
\end{enumerate}
Functorial descent along $\Mirror^6\to\Mirror^7\to\Mirror^8$ therefore
identifies a single $F_4$-orbit of soldered blocks across exceptional rungs.
\end{proposition}

\begin{proof}[Proof sketch]
\emph{(i)} Freudenthal closure (Lemma~\ref{lem:freudenthal}) embeds
$G_2\times F_4\subset\E{8}$; $G_2$ fixes $\Quat\subset\Oct$ and hence
$H_2(\Quat)$ under the Jordan envelope. \emph{(ii)} Product mirror acts on
$\E{8}\times\E{8}$ by factor exchange; the diagonal $\jalg$ chart and
$H_2(\Quat)$ block are $\sigma$-fixed, so soldering is unchanged on
$\hat Q$. \emph{(iii)} Cartan descent preserves the $1|10|16$ partition and
$F_4$ action on $\jalg$ (Theorem~\ref{thm:e6}); the $4$ soldering directions
split as $\mathfrak{m}_{+1}$ in Lemma~\ref{lem:goldstone-sigma}. Observer
index $\iota$ selects the representative within the $G_2$-conjugacy class
(Lemma~\ref{lem:quat-F4}).
\end{proof}

\begin{proposition}[Coset metric]\label{prop:coset-metric}
$K|_{\mathfrak{m}}$ is positive-definite and $H$-invariant. Lorentzian
signature is carried by $\sigma_{\mathrm{sol}}$, not by Wick rotation of $K$.
\end{proposition}

\begin{proposition}[No-go: coset metric alone]\label{prop:coset-nogo}
Mirror stabilisation at $\jalg$ and $\E{6}$ fixes the \emph{Euclidean}
coset metric $K|_{\mathfrak{m}}$ (Proposition~\ref{prop:coset-metric}) and
the $32=28+4$ mass split (Theorem~\ref{thm:planck}), but it does \emph{not}
select $d=4$, a Lorentzian soldering form, or the
$\mathrm{SL}(2,\Oct)\cong\Spin(1,9)$ frame. The $H_2(\Oct)\cong\R^{1,9}$
embedding of Definition~\ref{def:spacetime:H2O} motivates soldering; it is
not a theorem from void and mirror alone until observer-admitted lifetime
charts supply the clock block (Theorem~\ref{thm:spacetime:soldering-derived}).
\end{proposition}

\begin{theorem}[Soldering from mirror climb and lifetime chart]
\label{thm:spacetime:soldering-derived}
The mirror climb selects a canonical four-dimensional Lorentzian subspace
$H_2(\Quat)\subset H_2(\Oct)$ (Definition~\ref{def:spacetime:H2O}). For each
observer index $\iota\in\mathcal{I}$, the admitted lifetime flow
\eqref{eq:lifetime-admitted} of Definition~\ref{def:lifetime} couples
$\tau_\iota$ to the soldering clock
$\partial_{\tau_\iota}\sigma_{\mathrm{sol}}|_{H_2(\Quat)}$. Define
$\sigma_{\mathrm{sol}}$ to identify
$\orderparam^0,\ldots,\orderparam^3$ with this block. Then $d=4$, the
induced vierbein (Proposition~\ref{prop:vierbein}) carries signature
$(-,+,+,+)$, and soldering is \textbf{derived} from void, mirror, and
observer collapse---not postulated.
\end{theorem}

\begin{proof}[Proof sketch]
\emph{Step 1 (quaternionic rung).} $\Mirror^3(\Void)=\Quat$ has
$\dim_{\R}\Quat=4$, the first Hurwitz rung with an associative non-real
division algebra (Theorem~\ref{thm:rung-table}).

\emph{Step 2 (Jordan embedding).} Lemma~\ref{lem:quat-F4} fixes
$\Quat_\iota\subset\Oct$ and $H_2(\Quat)\subset H_2(\Oct)$ once the
observer index $\iota$ is chosen; mirror involution at $n=4\to5$ carries
the rung-$3$ datum $\Mirror^3(\Void)=\Quat$ into this block.

\emph{Step 3 (signature).} The restriction of $\det$ on $H_2(\Quat)$ is the
Minkowski norm; the four imaginary quaternion units supply the clock
directions. The coset metric $K|_{\mathfrak{m}}$ remains Euclidean;
Lorentzian signature enters through $\sigma_{\mathrm{sol}}$.

\emph{Step 4 (observer lifetime).} The index--lifetime map
(Theorem~\ref{thm:index-lifetime}) requires an admitted chart
$(\tau_\iota,g_\iota)$ whose flow \eqref{eq:lifetime-admitted} includes the
soldering clock term in $H_2(\Quat)$. Observer collapse
(Definition~\ref{def:observer}) on the $\sigma$-anti-invariant sector fixes
the initial datum $\orderparam_\star|_\iota$ and thereby selects the soldered
bundle $\mathcal{S}_\iota$ over which $g_\iota$ evolves
(Section~\ref{sec:observer:lifetime}).

\emph{Remaining gap.} Functorial transport across $n=6,7,8$ is settled by
Proposition~\ref{prop:soldering-transport}; global uniqueness reduces to the
tree-level massless $\mathfrak{m}_{+1}$ sector
(Theorem~\ref{thm:soldering-unique-reduced},
Conjecture~\ref{conj:soldering-unique}). The $\iota$-local embedding is
settled by Lemma~\ref{lem:quat-F4}.
\end{proof}

\begin{proposition}[Unique timelike soldering direction in $H_2(\Quat)$]
\label{prop:soldering-clock}
Fix $\iota\in\mathcal{I}$ and an admitted lifetime chart
$(\tau_\iota,g_\iota)$ of Definition~\ref{def:lifetime}. Along $g_\iota$,
the soldering time parameter $\tau_\iota$ couples to the zeroth clock
component $\orderparam^0$ via
\begin{equation}\label{eq:soldering-clock}
  \frac{d}{d\tau_\iota}\,\orderparam^0\big|_{g_\iota(\tau_\iota)}
  \;=\;
  \partial_{\tau_\iota}\sigma_{\mathrm{sol}}\big|_{H_2(\Quat)}.
\end{equation}
In the Lorentzian block $H_2(\Quat)$ of
Definition~\ref{def:spacetime:H2O}, the zeroth clock component
$\orderparam^0$ is the \emph{unique} timelike soldering direction admitted
by observer collapse: the $\sigma$-anti-invariant sector of $\hat Q$
breaks the $\Spin(1,9)$ frame to the $\iota$-restricted diagonal cell, and
capacity flow on $\mathcal{G}_{\mathrm{Albert}}$ fixes the sign convention
of $\det|_{H_2(\Quat)}$. The spatial components
$\orderparam^1,\ldots,\orderparam^3$ are the remaining spacelike directions
of the same block and do not carry the Page--Wootters clock.
\end{proposition}

\begin{proof}[Proof sketch]
Equation~\eqref{eq:soldering-clock} is the $\orderparam^0$-component of
\eqref{eq:lifetime-admitted}. On $H_2(\Quat)$, $\det$ has signature
$(-,+,+,+)$, so the timelike cone is one-dimensional up to sign;
Lemma~\ref{lem:quat-F4} and observer index $\iota$ fix the
$F_4$-equivariant orientation, and the above--below correspondence
(Definition~\ref{def:above-below}) aligns the negative eigen-direction
with the $\iota$-cell clock. No other component of
$\mathrm{Im}(\sigma_{\mathrm{sol}})$ is timelike.
\end{proof}

\begin{proposition}[Emergent vierbein]\label{prop:vierbein}
For a coset section $g(x)\in\E{6}$ with $g^{-1}\partial_\mu g=A_\mu+e_\mu$,
the $\mathfrak{m}$-component defines a composite vierbein. Under
Theorem~\ref{thm:spacetime:soldering-derived},
$\eta_{\mu\nu}=e^a_\mu e^b_\nu\eta_{ab}$ has signature $(-,+,+,+)$.
\end{proposition}

\begin{conjecture}[Functorial uniqueness of soldering]
\label{conj:soldering-unique}
Given $\iota\in\mathcal{I}$ and $\orderparam_\star|_\iota$, the soldering
block $H_2(\Quat)\subset H_2(\Oct)$, the clock coupling
\eqref{eq:soldering-clock}, and the timelike axis $\orderparam^0$ are
\textbf{already unique} by Lemma~\ref{lem:quat-F4} and
Proposition~\ref{prop:soldering-clock}. Functorial transport on
$n=6,7,8$ is supplied by Proposition~\ref{prop:soldering-transport}; the
residual conjecture is alignment of the tree-level massless
$\mathfrak{m}_{+1}$ Goldstones with the observer vierbein of
Proposition~\ref{prop:vierbein} and with $\iota$-sector lifetime charts
(Conjecture~\ref{conj:lifetime-unique}).
\end{conjecture}

\begin{theorem}[Soldering uniqueness, conditional reduction]
\label{thm:soldering-unique-reduced}
Assume Proposition~\ref{prop:soldering-transport},
Lemma~\ref{lem:quat-F4},
Proposition~\ref{prop:soldering-clock},
Theorem~\ref{thm:spacetime:soldering-derived}, and
Theorem~\ref{thm:lifetime-unique-reduced}. Then
Conjecture~\ref{conj:soldering-unique} reduces to the tree-level massless
sector only: the soldered block $H_2(\Quat)$, clock coupling
\eqref{eq:soldering-clock}, and $F_4$-equivariant vierbein of
Proposition~\ref{prop:vierbein} are unique on each $\iota$-sector once the
four $\mathfrak{m}_{+1}$ Goldstone directions of
Lemma~\ref{lem:goldstone-sigma} are included. Functorial transport on
$n=6,7,8$ is proved by Proposition~\ref{prop:soldering-transport}; the
residual open step is nonperturbative alignment of the massless
$\sigma=+1$ reservoir with the observer clock along $g_\iota$.
\end{theorem}

\begin{proof}[Proof sketch]
Proposition~\ref{prop:soldering-transport} closes transport of
$\sigma_{\mathrm{sol}}$ under automorphism, product, and Cartan mirrors.
Lemma~\ref{lem:quat-F4} and Proposition~\ref{prop:soldering-clock} fix the
$\iota$-local block and timelike axis. The vierbein of
Proposition~\ref{prop:vierbein} is determined on the massive $28$-direction
complement by Theorem~\ref{thm:rp}; Lemma~\ref{lem:goldstone-sigma}
isolates the residual $4$-dimensional $\mathfrak{m}_{+1}$ sector.
Theorem~\ref{thm:lifetime-unique-reduced} matches clock reparametrizations
along $g_\iota$. Any remaining ambiguity is confined to massless
Goldstone--vierbein alignment in the $\sigma=+1$ reservoir.
\end{proof}

\begin{remark}[Residual gap for the full soldering theorem]
\label{rem:soldering-full-gap}
Lemma~\ref{lem:quat-F4}, Proposition~\ref{prop:soldering-clock},
Proposition~\ref{prop:soldering-transport}, and
Theorem~\ref{thm:spacetime:soldering-derived} supply $d=4$, Lorentzian
signature, $\iota$-local clock coupling, and exceptional-rung transport
without postulate. What remains for a \emph{complete} functorial theorem is:
(i)~identify the resulting vierbein with the composite of
Proposition~\ref{prop:vierbein} on $\E{6}/H$ in the tree-level massless
$\mathfrak{m}_{+1}$ sector; and (ii)~close
Conjecture~\ref{conj:lifetime-unique} on the massless Goldstone reservoir.
Theorem~\ref{thm:soldering-unique-reduced} records the reduction; Open
Problem~\textbf{O2} is thereby confined to the massless sector.
\end{remark}

\begin{remark}[Epistemic status of P1]
Proposition~\ref{prop:coset-nogo} records the honest no-go for the coset
metric alone. Theorem~\ref{thm:spacetime:soldering-derived} \textbf{derives}
soldering and $d=4$ from mirror climb plus observer lifetime charts;
$\iota$-local uniqueness is supplied by Lemma~\ref{lem:quat-F4} and
Proposition~\ref{prop:soldering-clock}; exceptional-rung transport is supplied
by Proposition~\ref{prop:soldering-transport}. The residual open input is
tree-level massless-sector closure
(Theorem~\ref{thm:soldering-unique-reduced},
Conjecture~\ref{conj:soldering-unique},
Remark~\ref{rem:soldering-full-gap}, Open Problem~\textbf{O2} in
Section~\ref{sec:conclusions}).
\end{remark}

\begin{theorem}[Semiclassical reflection positivity along observer lifetimes]
\label{thm:rp}
Fix $\iota\in\mathcal{I}$ and a lifetime chart
$(\tau_\iota,g_\iota)$ admitted by capacity flow
(Definitions~\ref{def:lifetime},~\ref{def:lifetime-index}). On EIII, with
radiative mass $\mu^2>0$ on $28$ coset scalars (Coleman--Weinberg), the
Euclidean $\sigma$-model \eqref{eq:action} satisfies reflection positivity
along the observer clock $\tau_\iota$: the soldering map identifies the
soldered time component $\orderparam^0$ with $\tau_\iota$ on the
$\iota$-restricted trajectory $g_\iota(\tau_\iota)\in\mathcal{S}_\iota$, and
the semiclassical OS reflection plane is taken at fixed $\tau_\iota=0$ in
that chart. Equivalently, RP holds along the Page--Wootters clock carried by
the quaternionic block $H_2(\Quat)$ once the observer index $\iota$ selects
the diagonal cell.
\end{theorem}

\begin{proof}[Proof sketch]
The $\iota$-restricted gradient flow \eqref{eq:lifetime-admitted} is a
semiclassical saddle expansion about $\orderparam_\star|_\iota$ on the massive
$28$-dimensional coset complement of the $4$ soldering directions. With
$\mu^2>0$, the quadratic part defines a reflection-positive Gaussian measure;
nonnegativity of the capacity-derived interaction potential
(Theorem~\ref{thm:vacuum}) and compactness on EIII transfer positivity to the
perturbed measure in the standard OS semiclassical argument
\cite{GlimmJaffe1987,OsterwalderSchrader1973}. Functoriality of
Theorem~\ref{thm:index-lifetime} identifies the global clock
$\orderparam^0$ with the observer parameter $\tau_\iota$ along $g_\iota$.
\end{proof}

\begin{remark}[Honest semiclassical scope]
Theorem~\ref{thm:rp} is \emph{derived} in the radiatively gapped EIII regime;
it does not follow from void and mirror alone. Theorem~\ref{thm:rp-gamma-neg}
extends this to the full observable factor $d\mu_{\Gamma_{-1}}$ once
\eqref{eq:gamma-split} holds and massless modes are confined to
$d\mu_{\Gamma_{+1}}$ (Lemma~\ref{lem:goldstone-sigma}); the sole residual
input is nonperturbative verification of the measure split
(Conjecture~\ref{conj:rp}).
\end{remark}

\begin{proposition}[$\sigma$-sector split of the Euclidean measure]
\label{prop:sigma-measure}
Let $\sigma$ denote the product-mirror involution on the EIII configuration
space $\mathcal{M}$ and on the composite gauge bundle, acting consistently with
Theorem~\ref{thm:duality}. The Euclidean path measure for \eqref{eq:action},
formally
\begin{equation}\label{eq:gamma-full}
  d\mu_\Gamma[\orderparam,A]
  \;\propto\;
  \exp\!\bigl(-S[\orderparam,A]\bigr)\,\mathcal{D}\orderparam\,\mathcal{D}A,
\end{equation}
admits a $\sigma$-sector decomposition
\begin{equation}\label{eq:gamma-split}
  d\mu_\Gamma
  \;=\;
  d\mu_{\Gamma_{-1}}\otimes d\mu_{\Gamma_{+1}},
\end{equation}
where $d\mu_{\Gamma_{\pm1}}$ is the restriction to $\sigma=\pm1$
fluctuations about $\orderparam_\star$ and
$\Pi_{\sigma=-1}\,d\mu_{\Gamma_{-1}}=d\mu_{\Gamma_{-1}}$.
Observable correlators are $\Pi_{\sigma=-1}$-matrix elements of
$d\mu_{\Gamma_{-1}}$ only; the $\sigma=+1$ factor is a mirror-shadow
reservoir (Theorem~\ref{thm:duality}).
\end{proposition}

\begin{lemma}[Goldstones in the $\sigma$-invariant coset sector]
\label{lem:goldstone-sigma}
At the tree-level EIII saddle of Theorem~\ref{thm:vacuum}, the $32$
Hessian-null coset directions split under $\sigma$ as
\begin{equation}\label{eq:goldstone-split}
  \mathfrak{m}
  \;=\;
  \mathfrak{m}_{-1}\oplus\mathfrak{m}_{+1},
  \qquad
  \dim\mathfrak{m}_{-1}=28,\quad \dim\mathfrak{m}_{+1}=4,
\end{equation}
with $\mathfrak{m}_{+1}$ the four massless Goldstones of the soldering block
$H_2(\Quat)$ and their $\sigma$-invariant mirror doubles on
$\overline{\mathbf{16}}_{+3}$. After Coleman--Weinberg gapping, the
$28$ massive directions lie in $\mathfrak{m}_{-1}$ and are covered by
Theorem~\ref{thm:rp}. Full reflection positivity along $\tau_\iota$ on the
\emph{observable} sector therefore reduces to RP for $d\mu_{\Gamma_{-1}}$
(Proposition~\ref{prop:sigma-measure}); the $\sigma=+1$ Goldstone reservoir
enters only through normalization of the product \eqref{eq:gamma-split}.
\end{lemma}

\begin{theorem}[Reflection positivity of $d\mu_{\Gamma_{-1}}$, conditional]
\label{thm:rp-gamma-neg}
Fix $\iota\in\mathcal{I}$ and an admitted lifetime chart
$(\tau_\iota,g_\iota)$. Assume:
\begin{enumerate}[label=(\roman*),leftmargin=*]
  \item the $\sigma$-sector decomposition \eqref{eq:gamma-split} of
    Proposition~\ref{prop:sigma-measure};
  \item all tree-level massless coset modes lie in
    $\mathfrak{m}_{+1}\subset\mathfrak{m}$ and hence in $d\mu_{\Gamma_{+1}}$
    (Lemma~\ref{lem:goldstone-sigma}).
\end{enumerate}
Then $d\mu_{\Gamma_{-1}}$ is reflection-positive along the observer clock
$\tau_\iota$ of Definition~\ref{def:lifetime}: the $28$ radiatively gapped
$\sigma=-1$ coset directions of Theorem~\ref{thm:rp} together with the
$\sigma=-1$ matter sector define a reflection-positive Euclidean measure on
the observable factor, with reflection plane at fixed $\tau_\iota=0$ in chart
$g_\iota$.
\end{theorem}

\begin{proof}[Proof sketch]
By (ii), massless fluctuations do not enter $d\mu_{\Gamma_{-1}}$; the
observable factor is the product of the $28$-dimensional Coleman--Weinberg
gapped complement of $\mathfrak{m}_{+1}$ and the $\Pi_{\sigma=-1}$-projected
gauge--matter bundle. Theorem~\ref{thm:rp} supplies reflection positivity on
the massive coset block; $\sigma$-equivariance of \eqref{eq:action} and
compactness on EIII extend the OS semiclassical argument to the full
$d\mu_{\Gamma_{-1}}$ once (i) isolates the observable tensor factor. The
$\sigma=+1$ Goldstone reservoir normalizes the product and is
quotient-non-observable (Theorem~\ref{thm:duality}).
\end{proof}

\begin{proposition}[OS reconstruction on the observer sector]
\label{prop:os-observer}
Suppose the hypotheses of Theorem~\ref{thm:rp-gamma-neg} and the standard
OS axioms (Euclidean invariance, clustering along $\tau_\iota$, symmetry)
hold for $d\mu_{\Gamma_{-1}}$. Then $\Pi_{\sigma=-1}$-projected Schwinger
functions reconstruct a unitary Lorentzian quantum field theory on the
observer Hilbert space: the Osterwalder--Schrader theorem
\cite{OsterwalderSchrader1973} applies to the observable factor, with
Minkowski signature imported through the soldering map
$\sigma_{\mathrm{sol}}$ (Theorem~\ref{thm:spacetime:soldering-derived}) and
Page--Wootters time $\tau_\iota$. Mirror-shadow correlators in
$d\mu_{\Gamma_{+1}}$ do not enter physical matrix elements
(Proposition~\ref{prop:sigma-measure}).
\end{proposition}

\begin{conjecture}[Nonperturbative reflection positivity]\label{conj:rp}
Proposition~\ref{prop:sigma-measure}, Lemma~\ref{lem:goldstone-sigma},
Theorem~\ref{thm:rp-gamma-neg}, and Proposition~\ref{prop:os-observer}
reduce full OS reconstruction to verifying \eqref{eq:gamma-split} for the
composite $\Spin(10)\times U(1)$ gauge measure: once the product
decomposition holds and all massless modes lie in $\mathfrak{m}_{+1}$
(Lemma~\ref{lem:goldstone-sigma}), reflection positivity of
$d\mu_{\Gamma_{-1}}$ follows \textbf{conditionally} and
Proposition~\ref{prop:os-observer} supplies unitary Lorentzian dynamics on
the observer sector. The residual gap is nonperturbative proof of
\eqref{eq:gamma-split} itself---not a separate RP postulate for the massless
sector, which is quotient-non-observable in $d\mu_{\Gamma_{+1}}$
(Theorem~\ref{thm:duality}).
\end{conjecture}

%----------------------------------------------------------------------%
\subsection{Gauge, anomalies, gravity}
\label{sec:physics:gauge}
%----------------------------------------------------------------------%

\begin{theorem}[Anomaly cancellation]\label{thm:anomaly}
For one embedded $\mathbf{27}$: cubic $16-80+64=0$; mixed coefficient $-1$
cancelled by Green--Schwarz~\cite{GreenSchwarz1984}. With $n_F=3$ replicas of
the full $\mathbf{27}$ arithmetic, anomalies cancel family-wise.
\end{theorem}

\begin{theorem}[$b_0=12$]\label{thm:beta}
With $C_A(\Spin(10))=8$, $T(\mathbf{16})=2$, $n_F=3$, $n_S=28$:
$b_0=(88-24-28)/3=12$.
\end{theorem}

\begin{theorem}[Induced Planck mass]\label{thm:planck}
Sakharov induction from $28$ massive coset scalars gives
\begin{equation}\label{eq:planck}
  \MPl^2
  \;=\;
  \frac{28}{6}\cdot\frac{\mu^2}{16\pi^2}
  \log\frac{\Lambda^2}{\mu^2}.
\end{equation}
Physical scale identification ($\mu$, $\MPl$ in GeV) is deferred to
Section~\ref{sec:scales} (derivation step \textbf{D3}).
\end{theorem}

\begin{definition}[Dimensionless Newton coupling]\label{def:kappa}
At RG scale $k$, define
\begin{equation}\label{eq:kappa-def}
  \kappa(k)\;\equiv\;\frac{k^2}{\MPl^2(k)},
\end{equation}
with $\MPl(k)$ the running scale induced by the Sakharov flow of
Theorem~\ref{thm:planck}.
\end{definition}

\begin{theorem}[Gravitational fixed point]\label{thm:kappa}
Assume: (i)~Einstein--Hilbert truncation with coupling $\MPl^2(k)$;
(ii)~one-loop contribution from the $28$ massive coset scalars of
Theorem~\ref{thm:planck}; (iii)~higher-curvature operators and graviton loops
neglected; (iv)~anomalous dimension
$\eta_{\MPl}=-(5/48\pi^2)\kappa$ in this truncation. Then
\begin{equation}\label{eq:betakappa}
  \beta_\kappa\;\equiv\;k\frac{d\kappa}{dk}
  \;=\;2\kappa-\frac{5}{48\pi^2}\,\kappa^2,
\end{equation}
with $\kappa_\star=96\pi^2/5$ and
$\partial\beta_\kappa/\partial\kappa|_{\kappa_\star}=-2$.
\end{theorem}

\begin{definition}[Cosmological coupling]\label{def:cc}
The heat-kernel coefficient $a_0$ of the coset--gravity complex produces an
Einstein--Hilbert vacuum term $\int\sqrt{|g|}\,\Lambda_{\mathrm{cc}}$.
Treat $\Lambda_{\mathrm{cc}}$ as an independent dimension-4 coupling in the
effective action, not fixed by $\mu$ or $\MPl$.
\end{definition}

\begin{lemma}[$a_0$ heat-kernel projection to $\sigma=+1$]
\label{lem:a0-sigma-plus}
Let $\mathcal{D}$ denote the second-order operator of the coset--gravity
heat-kernel complex in the Sakharov truncation of Theorem~\ref{thm:planck},
acting on $\hat Q$. Mirror collapse
$\Pi_{\sigma=\pm1}$ (Definition~\ref{def:heterotic-collapse}) decomposes the
$a_0$ coefficient of $\mathrm{Tr}\,e^{-t\mathcal{D}}$ into
\begin{equation}\label{eq:a0-split}
  a_0(\mathcal{D})
  \;=\;
  a_0^{(+1)}+a_0^{(-1)},
\end{equation}
with $a_0^{(\pm1)}$ supported on $H^\bullet(\hat Q)^{\sigma=\pm1}$
(Theorem~\ref{thm:duality}). The induced vacuum-energy coupling in
Definition~\ref{def:cc} is the $\sigma$-labelled density
\begin{equation}\label{eq:a0-density}
  \Lambda_{\mathrm{cc}}(x)
  \;=\;
  a_0^{(+1)}(x)\;+\;a_0^{(-1)}(x),
\end{equation}
and the \emph{sequestered} effective action at $\kappa_\star$ is
\begin{equation}\label{eq:gamma-cc-eff}
  \Gamma_{\mathrm{eff}}
  \;=\;
  \Gamma_{\mathrm{coset}}
  \;+\;
  \int d^4x\,\sqrt{|g|}\,
  \Pi_{\sigma=+1}\,a_0^{(+1)}(x),
\end{equation}
so that $\Pi_{\sigma=-1}\Lambda_{\mathrm{cc}}=0$ on observable configurations
(Definition~\ref{def:physical}). Equation~\eqref{eq:gamma-cc-eff} is the
effective-action target for Conjecture~\ref{conj:cc-sequester}; it advances
open problem \textbf{O6} from a coupling hypothesis to a concrete projection
step on the heat-kernel density.
\end{lemma}

\begin{proof}[Proof sketch]
The coset--gravity complex is $\sigma$-equivariant under product mirror on
$\hat Q$. The $a_0$ coefficient is the $t^0$ term in the heat-kernel
expansion and inherits the $\mathbb{Z}_2$ decomposition because
$\mathrm{Tr}\,e^{-t\mathcal{D}}$ is a $\sigma$-invariant functional.
Observable fermions and scalars lie in $H^\bullet(\hat Q)^{\sigma=-1}$
(Definition~\ref{def:physical}); mirror partners in
$H^\bullet(\hat Q)^{\sigma=+1}$ are quotient-non-observable
(Theorem~\ref{thm:duality}). Routing $\Lambda_{\mathrm{cc}}$ through
$\Pi_{\sigma=+1}$ in \eqref{eq:gamma-cc-eff} removes the $a_0$ vacuum term
from $d\mu_{\Gamma_{-1}}$ while preserving normalization in
$d\mu_{\Gamma_{+1}}$ (Proposition~\ref{prop:sigma-measure}).
\end{proof}

\begin{theorem}[Vacuum energy RG relevance]\label{thm:cc-rg}
In the same truncation as Theorem~\ref{thm:planck}, the $a_0$ term is
scheme-dependent and does not cancel against the massive-mode $a_2$ sum. Under
canonical scaling,
\begin{equation}\label{eq:beta-Lambda}
  \beta_{\Lambda_{\mathrm{cc}}}
  =-4\,\Lambda_{\mathrm{cc}}+\mathcal{O}(\text{loops}),
\end{equation}
so $\Lambda_{\mathrm{cc}}$ is a \emph{relevant} direction: a UV fixed point
requires $\Lambda_{\mathrm{cc},\star}$ as input, not as output of void and
mirror alone.
\end{theorem}

\begin{proposition}[No-go: CC not from void and mirror alone]
\label{prop:cc-nogo}
The cosmological constant problem---explaining
$\Lambda_{\mathrm{obs}}\sim(10^{-3}\,\mathrm{eV})^4$ without fine-tuning---
cannot be resolved from the void and mirror axioms and derived $\E{6}$
dynamics \emph{alone}. Observable configurations lie in
$H^\bullet(\hat Q,V_{16})^{\sigma=-1}$ (Definition~\ref{def:physical}); the
heat-kernel $a_0$ vacuum energy and Theorem~\ref{thm:cc-rg} act on the full
quotient $\hat Q$, not only on the anti-invariant sector. Without a
sequestering mechanism that routes $\Lambda_{\mathrm{cc}}$ into the
non-observable $\sigma$-invariant sector (Theorem~\ref{thm:duality}),
$\Lambda_{\mathrm{cc}}$ remains a tuned environmental coupling in the
observer sector.
\end{proposition}

\begin{remark}[Observer-local FRG]\label{rmk:observer-frg}
Mirror collapse $\Pi_{\sigma=-1}$ (Definition~\ref{def:heterotic-collapse})
splits the effective action into $\sigma=\pm1$ sectors. An \emph{observer-local}
FRG is the flow restricted to anti-invariant modes along a lifetime chart
$g_\iota(\tau_\iota)$ (Definition~\ref{def:lifetime}). $\sigma$-invariant
(mirror-partner) modes are non-observable but may back-react on marginal
couplings---notably $\Lambda_{\mathrm{cc}}$---without entering the observable
Hilbert space.
\end{remark}

\begin{proposition}[$\Lambda_{\mathrm{cc}}$ decoupling in the $\sigma=+1$ sector]
\label{prop:cc-decouple}
Suppose the cosmological coupling $\Lambda_{\mathrm{cc}}$ couples only to
$\sigma$-invariant modes in the heat-kernel $a_0$ sector of
Definition~\ref{def:cc}---equivalently, its insertions in the effective
action are $\Pi_{\sigma=+1}$-projected along the observer-local FRG of
Remark~\ref{rmk:observer-frg}. Then along the $\sigma=-1$ observer flow
$g_\iota(\tau_\iota)$,
\begin{equation}\label{eq:beta-Lambda-decouple}
  \beta_{\Lambda_{\mathrm{cc}}}\Big|_{\sigma=-1}=0,
\end{equation}
and $\Lambda_{\mathrm{cc}}$ is not a relevant coupling in the observable
Jacobian row: the $\Lambda_{\mathrm{cc}}$ entry of $J$ vanishes on
$H^1(\hat Q,V_{16})^{\sigma=-1}$. Vacuum energy RG flows entirely within
$d\mu_{\Gamma_{+1}}$ (Proposition~\ref{prop:sigma-measure}), consistent with
Proposition~\ref{prop:cc-nogo}.
\end{proposition}

\begin{theorem}[$a_0$ sequestering at $\kappa_\star$, sketch]
\label{thm:cc-sequester-sketch}
At the gravitational fixed point $\kappa_\star=96\pi^2/5$
(Theorem~\ref{thm:kappa}), assume Lemma~\ref{lem:a0-sigma-plus} and the
measure split \eqref{eq:gamma-split}. Then the effective vacuum-energy
coupling in \eqref{eq:gamma-cc-eff} satisfies
\begin{equation}\label{eq:a0-sequester}
  \Lambda_{\mathrm{cc}}
  \;=\;
  \Pi_{\sigma=+1}\,\Lambda_{\mathrm{cc}},
\end{equation}
realizing the hypothesis of Proposition~\ref{prop:cc-decouple}:
$\Lambda_{\mathrm{cc}}$ couples only to $\sigma$-invariant modes and its RG
flow is confined to $d\mu_{\Gamma_{+1}}$. The observer sector sees
$\Pi_{\sigma=-1}\Lambda_{\mathrm{cc}}=0$ at the fixed point without tuning
$H^1(\hat Q,V_{16})^{\sigma=-1}$.
\end{theorem}

\begin{proof}[Proof sketch]
Mirror collapse $\Pi_{\sigma=-1}$
(Definition~\ref{def:heterotic-collapse}) splits the coset--gravity
heat-kernel complex into $\sigma=\pm1$ sectors on $\hat Q$
(Theorem~\ref{thm:duality}). At $\kappa_\star$, the one-loop Jacobian of
Proposition~\ref{prop:frg-jacobian} has $\Lambda_{\mathrm{cc},\star}=0$ as a
free Gaussian label; routing the $a_0$ density into
$\mathfrak{m}_{+1}\oplus\sigma$-invariant graviton modes leaves the
anti-invariant row untouched (Remark~\ref{rmk:observer-frg}).
Proposition~\ref{prop:cc-decouple} then yields
\eqref{eq:beta-Lambda-decouple}. The residual step is to derive
\eqref{eq:a0-sequester} from the nonperturbative effective action rather than
from the FRG coupling assignment.
\end{proof}

\begin{corollary}[Vanishing $\Lambda_{\mathrm{cc}}$ $\beta$-function on observable flow at FP]
\label{cor:beta-Lambda-fp}
At the extended fixed point $(\xi_\star,0,0,\kappa_\star,0)$ of
\eqref{eq:frg-fp}, under the $\Pi_{\sigma=+1}$ coupling hypothesis of
Proposition~\ref{prop:cc-decouple}---equivalently, under the $a_0$
sequestering route of Theorem~\ref{thm:cc-sequester-sketch}---the
observer-local FRG along $g_\iota(\tau_\iota)$ satisfies
\begin{equation}\label{eq:beta-Lambda-obs-fp}
  \beta_{\Lambda_{\mathrm{cc}}}\Big|_{\sigma=-1,\,\mathrm{FP}}
  \;=\;
  0,
\end{equation}
and $\Lambda_{\mathrm{cc}}$ is \emph{not} a UV-relevant direction in the
observable stability matrix: the $J$-row on
$H^1(\hat Q,V_{16})^{\sigma=-1}$ has no $\Lambda_{\mathrm{cc}}$ entry. The
global $\beta_{\Lambda_{\mathrm{cc}}}=-4\Lambda_{\mathrm{cc}}$ of
Theorem~\ref{thm:cc-rg} flows entirely in $d\mu_{\Gamma_{+1}}$.
\end{corollary}

\begin{proof}
Immediate from Proposition~\ref{prop:cc-decouple} at
$\Lambda_{\mathrm{cc},\star}=0$ and from the diagonal Jacobian
\eqref{eq:jacobian-diag} once the $\Lambda_{\mathrm{cc}}$ coupling is
$\sigma=+1$-projected along the observer flow
(Appendix~\ref{sec:appendix:frg}).
\end{proof}

\begin{conjecture}[CC sequestering at $\kappa_\star$]\label{conj:cc-sequester}
Lemma~\ref{lem:a0-sigma-plus}, Theorem~\ref{thm:cc-sequester-sketch}, and
Corollary~\ref{cor:beta-Lambda-fp} show that $\Pi_{\sigma=+1}$ projection of
the $a_0$ heat-kernel density at $\kappa_\star=96\pi^2/5$
(Theorem~\ref{thm:kappa}) would realize $\Lambda_{\mathrm{obs}}\approx0$
without tuning the anti-invariant sector (Theorem~\ref{thm:duality}).
Proposition~\ref{prop:cc-decouple} is thereby \emph{promoted from hypothesis
to target theorem}. The residual gap is a single effective-action step:
derive \eqref{eq:gamma-cc-eff}---equivalently \eqref{eq:a0-sequester}---from
the nonperturbative $\Gamma_{\mathrm{eff}}$ at $\kappa_\star$ by proving that
the $a_0^{(-1)}$ contribution cancels in $\Pi_{\sigma=-1}$-projected
observables while $a_0^{(+1)}$ is stored in $d\mu_{\Gamma_{+1}}$.
\end{conjecture}

Collect dimensionless couplings
$(\xi,\hat g^2,\hat y^2,\kappa,\Lambda_{\mathrm{cc}})$ with
$\xi=k^2/f^2$, $\hat g^2=k^2 g^2/f^2$, $\hat y^2=k^2 y^2/f^2$,
$\kappa$ as in \eqref{eq:kappa-def}, and $\Lambda_{\mathrm{cc}}$ as in
Definition~\ref{def:cc}. In the extended one-loop coset--gauge--gravity
truncation (still neglecting $R^2$, graviton loops, and higher curvature),
\begin{equation}\label{eq:frg-system}
\begin{aligned}
  \beta_{\xi}&=2\xi,
  &\quad
  \beta_{\hat g^2}&=2\hat g^2+\frac{b_0}{48\pi^2}\hat g^4,
  &\quad
  \beta_{\hat y^2}&=2\hat y^2+\mathcal{O}(\hat y^4,\hat g^2\hat y^2),
  \\[4pt]
  \beta_{\kappa}&=2\kappa-\frac{5}{48\pi^2}\kappa^2,
  &\quad
  \beta_{\Lambda_{\mathrm{cc}}}&=-4\,\Lambda_{\mathrm{cc}}+\mathcal{O}(\text{loops}).
\end{aligned}
\end{equation}

\begin{proposition}[Jacobian in extended one-loop truncation]
\label{prop:frg-jacobian}
At the fixed point
$(\xi_\star,\hat g^2_\star,\hat y^2_\star,\kappa_\star,\Lambda_{\mathrm{cc},\star})
=(\xi_\star,0,0,96\pi^2/5,0)$, the stability matrix
$J=\partial\beta_i/\partial g_j$ has diagonal spectrum
$\mathrm{eig}(J)=\{+2,+2,+2,-2,-4\}$. With
$\theta_i=-\mathrm{eig}(\partial\beta_i/\partial g_i)$, exactly two directions
are UV-relevant: $\kappa$ ($\theta_\kappa=+2$) and
$\Lambda_{\mathrm{cc}}$ ($\theta_{\Lambda}=+4$). Details in
Appendix~\ref{sec:appendix:frg}.
\end{proposition}

\begin{conjecture}[Graviton/$R^2$ persistence of $\kappa_\star$]
\label{conj:kappa}
Proposition~\ref{prop:kappa-R2} and Corollary~\ref{cor:beta-Lambda-fp}
establish leading-order stability of $\kappa_\star=96\pi^2/5$ with
$\partial\beta_\kappa/\partial\kappa|_{\kappa_\star}=-2$ once
$\sigma$-invariant graviton and $R^2$ modes decouple from the observer-local
flow on $H^1(\hat Q,V_{16})^{\sigma=-1}$
(Remark~\ref{rmk:observer-frg}). Theorem~\ref{thm:cc-sequester-sketch} routes
subleading vacuum-energy back-reaction into $d\mu_{\Gamma_{+1}}$ without
modifying $\beta_\kappa|_{\sigma=-1}$ at the fixed point. The residual gap is
nonperturbative curved-background FRG confirmation that
$\beta_\kappa|_{\kappa_\star,\,\sigma=-1}=0$ in the full graviton/$R^2$
extension. Because $\kappa_\star\approx189\gg1$, such persistence is not a
consequence of void and mirror; it requires the extended FRG input of
Appendix~\ref{sec:appendix:frg}.
\end{conjecture}

%----------------------------------------------------------------------%
\subsection{GUT breaking from observer collapse}
\label{sec:physics:breaking}
%----------------------------------------------------------------------%

The heterotic product mirror supplies exact
$\mathbf{16}/\overline{\mathbf{16}}$ duality at $\Mirror^7(\Void)$; observer
collapse $\Pi_{\sigma=-1}$ (Definition~\ref{def:heterotic-collapse}) breaks
this pairing and selects chiral descent directions. Symmetry breaking below is
gradient flow of $V$ on $\iota$-restricted coset directions, functorially
descended from the global EIII minimum above
(Definition~\ref{def:above-below}).

\begin{lemma}[Hessian sector decomposition]\label{lem:breaking-hess}
Let $\orderparam_\star\in\mathcal{M}$ be the rank-one EIII vacuum
(Theorem~\ref{thm:vacuum}) and write fluctuations
$\delta\orderparam=\delta\orderparam_{\mathbf{1}}+\delta\orderparam_{\mathbf{10}}+\delta\orderparam_{\mathbf{16}}$
under the capacity partition
$\mathbf{27}=\mathbf{1}\oplus\mathbf{10}\oplus\mathbf{16}$
(Theorem~\ref{thm:partition}). On the coset
$\mathfrak{m}=\mathbf{16}_{-3}\oplus\overline{\mathbf{16}}_{+3}$, the Hessian
$\Hess V(\orderparam_\star)$ decomposes $H$-equivariantly into $\mathbf{1}$,
$\mathbf{10}$, and $\mathbf{16}\oplus\overline{\mathbf{16}}$ blocks. Restricting
to the observer diagonal $\mathfrak{d}=\mathrm{span}\{e_1,e_2,e_3\}$
(Definition~\ref{def:observer-index}), one obtains a further $\iota$-sector
split
\begin{equation}\label{eq:breaking-hess-iota}
  \Hess V(\orderparam_\star)\big|_{\mathfrak{d}}
  \;=\;
  \bigoplus_{\iota=1}^3 \Hess_\iota V(\orderparam_\star),
\end{equation}
with $\Hess_\iota$ acting on the $\iota$-cell of
$H^1(\hat Q,V_{16})^{\sigma=-1,\,\iota}$. At tree level
$\Hess V(\orderparam_\star)=0$ on $\mathfrak{m}$: the $32$ coset directions are
exact Goldstones; unstable modes enter only after Coleman--Weinberg radiative
correction~\cite{ColemanWeinberg1973}.
\end{lemma}

\begin{proposition}[Chiral breaking directions]\label{prop:breaking-chiral}
The observer projector $\Pi_{\sigma=-1}$ selects breaking along
$\sigma$-anti-invariant coset directions in
$H^1(\hat Q,V_{16})^{\sigma=-1}$ and suppresses the mirror-symmetric
$\overline{\mathbf{16}}$ partner. A vacuum-expectation trajectory $g_\iota(\tau)$
solving \eqref{eq:lifetime-admitted} therefore breaks
$H=\Spin(10)\times U(1)$ along chiral $\mathbf{16}$-type fluctuations tied to
slot $\iota$, not along a simultaneous
$\mathbf{16}\oplus\overline{\mathbf{16}}$ pairing. Duality is exact at the
heterotic rung (Theorem~\ref{thm:duality}); physical breaking is the localized
violation induced by collapse.
\end{proposition}

\begin{proof}[Proof sketch]
$\Pi_{\sigma=-1}=\tfrac12(\mathrm{id}-\sigma)$ annihilates $\sigma$-invariant
modes. On $\mathfrak{m}$ the mirror involution exchanges
$\mathbf{16}_{-3}$ and $\overline{\mathbf{16}}_{+3}$
(Proposition~\ref{prop:mirror:involution}); anti-invariant fluctuations lie in
$\mathbf{16}_{-3}$ modulo the $\iota$-diagonal restriction. The lifetime flow
\eqref{eq:lifetime-admitted} restricts $\nabla V$ to the $\iota$-sector,
breaking product mirror symmetry while preserving capacity balance on the global
datum $\mathfrak{U}^\star$.
\end{proof}

\begin{theorem}[First breaking, conditional]\label{thm:breaking-e6so10}
Assume the EIII vacuum $\orderparam_\star$ of Theorem~\ref{thm:vacuum}, observer
collapse on $\hat Q$, and a fixed observer index $\iota\in\mathcal{I}$. If
radiative corrections lift the $\mathbf{16}$-sector Hessian with a single
unstable mode per $\iota$-cell, then gradient flow of $V$ on the
$\iota$-restricted coset triggers the first step
\begin{equation}\label{eq:breaking-e6so10}
  \E{6}\;\longrightarrow\; \Spin(10)\times U(1)
\end{equation}
along the rank-one orbit of $\mathcal{M}$, with $U(1)$ charge aligned to the
$\iota$-diagonal of $\mathfrak{d}$. The residual stabilizer
$\mathfrak{h}=\mathfrak{so}_{10}\oplus\mathfrak{u}_1$ is the symmetry group of
Theorem~\ref{thm:universe}(iii).
\end{theorem}

\begin{proof}[Proof sketch]
Rank-one EIII fixes $\mathcal{M}=\E{6}/(\Spin(10)\times U(1))$ globally
(Theorem~\ref{thm:vacuum}). An $\iota$-localized unstable fluctuation in
$\mathbf{16}_{-3}$ is an adjoint-direction Goldstone of $\E{6}/H$; flow along
that mode restores a $\Spin(10)$-invariant submanifold while breaking the
remaining $U(1)_\iota\subset U(1)$ tied to the chosen diagonal. The branching
matches the standard
$\mathbf{78}=\mathbf{45}\oplus\mathbf{1}\oplus\mathbf{16}\oplus\overline{\mathbf{16}}$
pattern at the $\Spin(10)$-invariant vacuum~\cite{Witten1985,GeorgiGlashow1974}.
\end{proof}

\begin{lemma}[Unstable Hessian directions under $\Spin(10)$ branching]
\label{lem:breaking-hess-spin10}
Let $V_{\mathrm{eff}}$ denote the Coleman--Weinberg effective potential about
$\orderparam_\star$ (Theorem~\ref{thm:rp}). Write Jordan fluctuations
$\delta\orderparam\in\mathbf{27}=\mathbf{1}\oplus\mathbf{10}\oplus\mathbf{16}$
(Theorem~\ref{thm:partition}) and restrict to the observer diagonal
$\mathfrak{d}$ as in \eqref{eq:breaking-hess-iota}. After radiative lifting,
each $\Hess_\iota V_{\mathrm{eff}}(\orderparam_\star)$ is
$H$-equivariant on its $\iota$-cell. Under the standard $\Spin(10)$ branching
\begin{equation}\label{eq:breaking-adjoint-branch}
  \mathbf{78}
  \;=\;
  \mathbf{45}\oplus\mathbf{1}\oplus\mathbf{16}\oplus\overline{\mathbf{16}},
  \qquad
  \mathbf{27}=\mathbf{1}\oplus\mathbf{10}\oplus\mathbf{16},
\end{equation}
and the observer projector $\Pi_{\sigma=-1}$, unstable directions classify as
follows:
\begin{enumerate}[label=(\roman*),leftmargin=*]
  \item $\iota=3$ ($\mathbf{16}$-sector, capacity weight $16/27$): unstable
    modes lie in $\mathbf{16}_{-3}$ and embed in the adjoint
    $\mathbf{16}\subset\mathbf{78}$; they couple to the $\mathbf{45}$ adjoint
    orbit of $\Spin(10)$.
  \item $\iota=2$ ($\mathbf{10}$-sector, weight $10/27$): unstable modes lie in
    $\mathbf{10}_{-2}\subset\mathbf{27}$ and embed in the antisymmetric
    $\mathbf{126}$ via $\mathbf{16}\times\mathbf{16}\supset\mathbf{126}$
    (with $\overline{\mathbf{16}}$ suppressed by $\Pi_{\sigma=-1}$).
  \item $\iota=1$ ($\mathbf{1}$-sector, weight $1/27$): unstable trace modes
    couple to the $\mathbf{10}$ fundamental through the Yukawa portal
    $\mathbf{16}\times\mathbf{16}\times\mathbf{10}$ and the residual
    $\mathbf{10}\subset\mathbf{27}$ block.
\end{enumerate}
Mirror partners $\overline{\mathbf{16}}_{+3}$ do not appear in the unstable
spectrum of the observable sector.
\end{lemma}

\begin{proof}[Proof sketch]
$V_{\mathrm{eff}}$ is $H$-invariant, so each $\Hess_\iota V_{\mathrm{eff}}$ is
an $H$-intertwiner on the $\iota$-cell. The capacity partition aligns
$\iota=1,2,3$ with the $\mathbf{1}$, $\mathbf{10}$, $\mathbf{16}$ blocks of
$\mathbf{27}$ (Proposition~\ref{prop:above-below}). Items (i)--(iii) follow
from \eqref{eq:breaking-adjoint-branch}, the standard $\Spin(10)$ tensor
products $\mathbf{16}\times\mathbf{16}=\mathbf{1}\oplus\mathbf{45}\oplus
\mathbf{126}$ and $\mathbf{16}\times\mathbf{10}=\mathbf{16}\oplus\mathbf{144}$,
and Proposition~\ref{prop:breaking-chiral}.
\end{proof}

\begin{proposition}[Higgs labels from $\mathbf{27}$-sector Hessians]
\label{prop:breaking-higgs-labels}
Under Lemma~\ref{lem:breaking-hess-spin10}, radiatively unstable
$\Hess_\iota V_{\mathrm{eff}}$ directions map to observable Higgs multiplets
of the $\Spin(10)$ chain:
\begin{equation}\label{eq:breaking-higgs-map}
  \Hess_3 V_{\mathrm{eff}}\;\longmapsto\;\Phi_{\mathbf{45}}\in\mathbf{45}_H,
  \qquad
  \Hess_2 V_{\mathrm{eff}}\;\longmapsto\;\Phi_{\mathbf{126}}\in\mathbf{126}_H,
  \qquad
  \Hess_1 V_{\mathrm{eff}}\;\longmapsto\;\Phi_{\mathbf{10}}\in\mathbf{10}_H.
\end{equation}
Each label is the $\Pi_{\sigma=-1}$-projected unstable mode of the
corresponding $\mathbf{27}$ block, with capacity weight
$w_\iota\in\{1/27,10/27,16/27\}$ fixing the relative scale of the induced
VEV along $g_\iota$.
\end{proposition}

\begin{proof}[Proof sketch]
The $\mathbf{45}_H$ identification uses the adjoint embedding
$\mathbf{16}\subset\mathbf{78}$: an $\iota=3$ fluctuation is a
$\Spin(10)$-breaking direction in $\mathfrak{so}_{10}$. The $\mathbf{126}_H$
identification uses $\mathbf{16}\times\mathbf{16}\supset\mathbf{126}$ on the
$\iota=2$ cell. The $\mathbf{10}_H$ identification follows from the residual
$\mathbf{10}$ block at $\iota=1$ and the cubic portal
$\mathbf{16}\times\mathbf{16}\times\mathbf{10}$. Capacity weights set the
ordering of instability thresholds along the lifetime chart.
\end{proof}

\begin{theorem}[Second breaking, conditional]\label{thm:breaking-so10su5}
Assume Theorem~\ref{thm:breaking-e6so10} and a $\Spin(10)\times U(1)$
intermediate vacuum. If Coleman--Weinberg lifting produces an unstable
$\Hess_3 V_{\mathrm{eff}}$ mode on the $\iota=3$ cell---the
$\mathbf{16}$-sector with largest capacity weight $16/27$---then gradient flow
of $V$ along $g_3$ triggers
\begin{equation}\label{eq:breaking-so10su5}
  \Spin(10)\times U(1)\;\longrightarrow\;\mathrm{SU}(5)\times U(1)_X,
\end{equation}
with Higgs fluctuation $\Phi_{\mathbf{45}}\in\mathbf{45}_H$ of
Proposition~\ref{prop:breaking-higgs-labels}. The residual gauge algebra is
$\mathfrak{su}_5\oplus\mathfrak{u}_1$.
\end{theorem}

\begin{proof}[Proof sketch]
The $\iota=3$ unstable mode lies in $\mathbf{16}_{-3}$ and embeds in
$\mathbf{45}\subset\mathbf{78}$ (Lemma~\ref{lem:breaking-hess-spin10}(i)).
Flow along $\Phi_{\mathbf{45}}$ aligns the vacuum with a $\mathrm{SU}(5)$
subalgebra of $\mathfrak{so}_{10}$; the standard branching
$\mathbf{45}\to\mathbf{24}\oplus\mathbf{1}\oplus\cdots$ under
$\mathrm{SU}(5)\times U(1)_X$ matches Georgi--Glashow
$\Spin(10)\to\mathrm{SU}(5)\times U(1)_X$~\cite{GeorgiGlashow1974}. The
$16/27$ weight dominates the $\iota=2$ and $\iota=1$ thresholds, so this is
the second step after \eqref{eq:breaking-e6so10}.
\end{proof}

\begin{lemma}[$\mathbf{126}_H$ decomposition under $\Pi_{\sigma=-1}$]
\label{lem:breaking-126-sigma}
Let $\Phi_{\mathbf{126}}\in\mathbf{126}_H$ be the antisymmetric Higgs direction
of Proposition~\ref{prop:breaking-higgs-labels}, realized as the unstable
$\Hess_2 V_{\mathrm{eff}}$ mode on the $\iota=2$ cell. The $\Spin(10)$ tensor
product $\mathbf{16}\times\mathbf{16}=\mathbf{1}\oplus\mathbf{45}\oplus
\mathbf{126}$ supplies $\mathbf{126}_H$ as the antisymmetric wedge of chiral
spinors. Under the standard $\mathrm{SU}(5)\times U(1)_X$ branching,
\begin{equation}\label{eq:breaking-126-branch}
  \mathbf{126}
  \;\longrightarrow\;
  \mathbf{10}_{-4}\oplus\overline{\mathbf{10}}_{4}
  \oplus\mathbf{15}_{-6}\oplus\overline{\mathbf{15}}_{6}
  \oplus\mathbf{1}_{10},
\end{equation}
with $\mathbf{1}_{10}$ the $B\!-\!L=2$ singlet that lowers rank. The observer
projector $\Pi_{\sigma=-1}$ acts on \eqref{eq:breaking-126-branch} as follows:
\begin{enumerate}[label=(\roman*),leftmargin=*]
  \item Bilinear fluctuations $\omega_i\wedge\omega_j$ with
    $|\omega_i\rangle,|\omega_j\rangle\in H^1(\hat Q,V_{16})^{\sigma=-1}$
    are $\sigma$-invariant, so the $\mathbf{126}$ portal is built entirely from
    anti-invariant observer modes; mirror partners in
    $H^1(\hat Q,V_{16})^{\sigma=+1}$ do not enter.
  \item Conjugate multiplets paired by $\sigma$ on $\mathfrak{m}$ are suppressed:
    $\Pi_{\sigma=-1}$ annihilates the $\overline{\mathbf{10}}_{4}$ and
    $\overline{\mathbf{15}}_{6}$ mirror channels, retaining the chiral
    $\mathbf{10}_{-4}\oplus\mathbf{15}_{-6}\oplus\mathbf{1}_{10}$ sector.
  \item The $\mathrm{SU}(5)\to$SM breaking direction is the
    $\Pi_{\sigma=-1}$-projected component of $\mathbf{1}_{10}$ together with
    the $(1,1)$ singlet inside $\mathbf{15}_{-6}$; the residual
    $\mathrm{SU}(3)_c\times\mathrm{SU}(2)_L\times U(1)_Y$ algebra is the
    stabilizer after its VEV.
\end{enumerate}
Capacity weight $w_2=10/27$ fixes the scale ordering of this mode relative to
the $\iota=3$ and $\iota=1$ thresholds.
\end{lemma}

\begin{proof}[Proof sketch]
Item~(i) uses $\sigma(\omega_i\wedge\omega_j)=\sigma(\omega_i)\sigma(\omega_j)
\omega_i\wedge\omega_j$ with $\sigma(\omega_i)=-1$ on physical fermions.
Item~(ii) follows from Proposition~\ref{prop:breaking-chiral} and the mirror
exchange on $\overline{\mathbf{16}}_{+3}$. Item~(iii) is the standard
$\mathrm{SU}(5)$ rank-breaking pattern: $\langle\mathbf{1}_{10}\rangle$ removes
$U(1)_X$ and the $(1,1)\subset\mathbf{15}$ direction completes
$\mathrm{SU}(5)\to\mathrm{SU}(3)_c\times\mathrm{SU}(2)_L\times U(1)_Y$
\cite{GeorgiGlashow1974,Witten1985}. The $10/27$ weight is the $\mathbf{10}$
block capacity datum of Theorem~\ref{thm:partition}.
\end{proof}

\begin{theorem}[Third breaking, conditional]\label{thm:breaking-su5sm}
Assume Theorem~\ref{thm:breaking-so10su5} and a
$\mathrm{SU}(5)\times U(1)_X$ intermediate vacuum. If Coleman--Weinberg lifting
produces an unstable $\Hess_2 V_{\mathrm{eff}}$ mode on the $\iota=2$ cell---the
$\mathbf{10}$-sector with capacity weight $10/27$---then gradient flow of $V$
along $g_2$ triggers
\begin{equation}\label{eq:breaking-su5sm}
  \mathrm{SU}(5)\times U(1)_X
  \;\longrightarrow\;
  \mathrm{SU}(3)_c\times\mathrm{SU}(2)_L\times U(1)_Y,
\end{equation}
with Higgs fluctuation $\Phi_{\mathbf{126}}\in\mathbf{126}_H$ of
Proposition~\ref{prop:breaking-higgs-labels} and observable content given by
Lemma~\ref{lem:breaking-126-sigma}. The residual gauge algebra is
$\mathfrak{su}_3\oplus\mathfrak{su}_2\oplus\mathfrak{u}_1$.
\end{theorem}

\begin{proof}[Proof sketch]
The $\iota=2$ unstable mode lies in $\mathbf{10}_{-2}\subset\mathbf{27}$ and
embeds in $\mathbf{126}$ via $\mathbf{16}\times\mathbf{16}$
(Lemma~\ref{lem:breaking-hess-spin10}(ii)). Flow along
$\Pi_{\sigma=-1}\Phi_{\mathbf{126}}$ aligns the vacuum with the
$\mathbf{1}_{10}\oplus(1,1)\subset\mathbf{15}$ direction of
\eqref{eq:breaking-126-branch}; the standard branching
$\mathbf{15}\to(3,1)_{-1/3}\oplus(1,3)_{1/3}\oplus(1,1)_{-4/3}$ and
$\mathbf{10}\to(3,1)_{1/3}\oplus(1,2)_{-1/2}\oplus(1,1)_{2/3}$ under
$\mathrm{SU}(3)_c\times\mathrm{SU}(2)_L\times U(1)_Y$ matches Georgi--Glashow
$\mathrm{SU}(5)\to$SM~\cite{GeorgiGlashow1974}. The $10/27$ weight places this
step after \eqref{eq:breaking-so10su5} and before the $\iota=1$
electroweak stage.
\end{proof}

\begin{proposition}[Weinberg angle at $M_{\mathrm{GUT}}$]
\label{prop:weinberg-gut}
Assume Theorems~\ref{thm:breaking-e6so10} and~\ref{thm:breaking-so10su5}, with
$\Spin(10)$ gauge coupling normalization surviving to the GUT scale
(Theorem~\ref{thm:beta}, $b_0=12$). At $M_{\mathrm{GUT}}$,
\begin{equation}\label{eq:weinberg-gut}
  \sin^2\theta_W(M_{\mathrm{GUT}})=\tfrac38,
\end{equation}
where $\theta_W$ is the electroweak mixing angle in the canonical
$\mathrm{SU}(5)\subset\Spin(10)$ embedding. This is the group-theoretic anchor
for Prediction~\ref{pred:gut}.
\end{proposition}

\begin{proof}[Proof sketch]
In the $\Spin(10)$ adjoint, embed
$\mathrm{SU}(5)\times U(1)_X\subset\Spin(10)$ with generators normalized by
$\mathrm{Tr}(T_aT_b)=\delta_{ab}$ on the spinor $\mathbf{16}$. Hypercharge
$Y$ is the standard $\mathrm{SU}(5)$ combination; electric charge
$Q=T_3+Y$. Then
$\sin^2\theta_W=\mathrm{Tr}(T_3^2)/\mathrm{Tr}(Q^2)=3/8$ at unification,
independent of the intermediate $\mathbf{126}_H$ VEV scale
\cite{GeorgiGlashow1974,Witten1985}.
\end{proof}

\begin{theorem}[Electroweak breaking, sketch]\label{thm:breaking-ew}
Assume Theorems~\ref{thm:breaking-so10su5} and~\ref{thm:breaking-su5sm}.
If radiative lifting yields an unstable $\Hess_1 V_{\mathrm{eff}}$ mode on the
$\iota=1$ cell, then gradient flow along $g_1$ with
$\Phi_{\mathbf{10}}\in\mathbf{10}_H$ triggers
\begin{equation}\label{eq:breaking-ew}
  \mathrm{SU}(2)_L\times U(1)_Y
  \;\longrightarrow\;
  U(1)_{\mathrm{em}},
\end{equation}
with $\mathrm{SU}(3)_c$ unbroken. The electroweak scale is set by the
$1/27$ capacity weight on the final $\iota=1$ breaking stage.
\end{theorem}

\begin{proof}[Proof sketch]
$\mathbf{10}_H$ contains the $\mathrm{SU}(5)$ Higgs doublet
$\mathbf{5}\oplus\overline{\mathbf{5}}$; an $\iota=1$ fluctuation in
$\Phi_{\mathbf{10}}$ breaks the electroweak subgroup while leaving
$\mathrm{SU}(3)_c$ intact. The ordering $\iota=3\to2\to1$ follows capacity
weights $(16/27,10/27,1/27)$; the electroweak step is last because the
$\mathbf{1}$-sector carries the smallest weight.
\end{proof}

\begin{conjecture}[Derived GUT--Higgs chain]\label{conj:breaking-derived}
Sequential $\iota$-restricted gradient flow on the $28$ radiatively massive
coset directions (Coleman--Weinberg, as in Theorem~\ref{thm:rp}) realizes the
staged chain
\begin{equation}\label{eq:breaking-chain}
  \E{6}\to\Spin(10)\times U(1)\to SU(5)\to SU(3)_c\times SU(2)_L\times U(1)_Y
  \to SU(3)_c\times U(1)_{\mathrm{em}}
\end{equation}
as follows:
\begin{enumerate}[label=\textbf{(\arabic*)},leftmargin=*]
  \item \textbf{$\E{6}\to\Spin(10)\times U(1)$} --- any $\iota$-cell unstable
    $\mathbf{16}$ mode (Theorem~\ref{thm:breaking-e6so10}).
  \item \textbf{$\Spin(10)\to\mathrm{SU}(5)\times U(1)_X$} --- $\iota=3$
    $\mathbf{45}_H$ direction, largest weight $16/27$
    (Theorem~\ref{thm:breaking-so10su5}).
  \item \textbf{$\mathrm{SU}(5)\to\mathrm{SU}(3)_c\times\mathrm{SU}(2)_L
    \times U(1)_Y$} --- $\iota=2$ $\mathbf{126}_H$ direction, weight $10/27$
    (Theorem~\ref{thm:breaking-su5sm},
    Lemma~\ref{lem:breaking-126-sigma}).
  \item \textbf{Electroweak} --- $\iota=1$ $\mathbf{10}_H$ direction, weight
    $1/27$ (Theorem~\ref{thm:breaking-ew}).
\end{enumerate}
Observable Higgs content $\mathbf{45}_H$, $\mathbf{126}_H$, $\mathbf{10}_H$ is
the $\Pi_{\sigma=-1}$-projected unstable spectrum of
\eqref{eq:breaking-higgs-map}. Yukawa matrices $y_{ij}$ arise from cubic
overlap integrals
\begin{equation}\label{eq:breaking-yukawa}
  y_{ij}\;\propto\;\int_{\hat Q}\omega_\iota\wedge\omega_j\wedge\Phi_{\mathrm{Higgs}},
\end{equation}
of observer modes $|\omega_\iota\rangle$ on $\hat Q$
(Definition~\ref{def:observer}) against $\Phi_{\mathrm{Higgs}}\in\{\Phi_{\mathbf{45}},
\Phi_{\mathbf{126}},\Phi_{\mathbf{10}}\}$. Section~\ref{sec:yukawa} develops
the leading-order evaluation of \eqref{eq:breaking-yukawa}; the
$\mathrm{SU}(5)\to$SM stage is supplied by Theorem~\ref{thm:breaking-su5sm}
(conditional on Coleman--Weinberg lifting). Proposition~\ref{prop:weinberg-gut}
fixes $\sin^2\theta_W=3/8$ at $M_{\mathrm{GUT}}$ once $\Spin(10)$ survives.
The residual conjectural content is sequential realization of all four steps by
a single $\iota$-ordered gradient flow.
\end{conjecture}

\begin{conjecture}[GUT breaking]\label{conj:breaking}
Gradient flow of $V$ on $\iota$-restricted coset directions triggers
\eqref{eq:breaking-chain} with Higgs content $\mathbf{45}_H$, $\mathbf{126}_H$,
$\mathbf{10}_H$; Conjecture~\ref{conj:breaking-derived} supplies the
observer-derived mechanism (Lemma~\ref{lem:breaking-hess},
Lemma~\ref{lem:breaking-hess-spin10},
Proposition~\ref{prop:breaking-chiral},
Proposition~\ref{prop:breaking-higgs-labels},
Theorems~\ref{thm:breaking-e6so10},~\ref{thm:breaking-so10su5},
~\ref{thm:breaking-su5sm},~\ref{thm:breaking-ew},
Proposition~\ref{prop:weinberg-gut}).
\end{conjecture}

\begin{remark}[As above, so below]\label{rem:breaking-above-below}
Global capacity balance fixes the EIII minimum $\orderparam_\star\in\mathcal{M}$
\emph{above} (Theorem~\ref{thm:vacuum}); observer collapse localizes breaking
\emph{below} to the $\iota$-sector lifetime chart
(Definition~\ref{def:lifetime}). The correspondence \eqref{eq:above-below}
identifies the profinite universe datum $\mathfrak{U}^\star$ with fixed-point
collapse on $\hat Q$: the same potential $V$ governs the global vacuum and the
slot-resolved descent trajectories $g_\iota$. What is proved globally on
$\mathfrak{U}^\star$---the $1|10|16$ split, rank-one EIII, three observer
slots---appears locally as chiral gauge breaking in slot $\iota$, without
importing mirror-symmetric partners.
\end{remark}

%----------------------------------------------------------------------%
\subsection{Flavon hierarchy from observer capacity flow}
\label{sec:physics:flavon}
%----------------------------------------------------------------------%

Assume Conjecture~\ref{conj:breaking}. The flavour group
$SU(3)_F\subset\Spin(10)$ acts on the three chiral $\mathbf{16}$s; each
observer index $\iota\in\mathcal{I}$ labels one diagonal direction of
$\mathfrak{d}\subset\jalg$ and hence one $SU(3)_F$ flavour axis.

\begin{proposition}[Flavon sector from observer indices]\label{prop:flavon-sector}
The three observer slots of Theorem~\ref{thm:chirality-derived} index three
$SU(3)_F$ breaking stages. Along the lifetime chart $g_\iota(\tau_\iota)$
(Definition~\ref{def:lifetime}), the capacity flow
$\omega^\star_\iota(\tau)=\capacity_\iota(\tau)$ of
Definition~\ref{def:lifetime-index} records which flavon direction is active at
time $\tau$ within slot $\iota$. Mirror stabilisation fixes $n_F=3$ but not,
by itself, Yukawa eigenvalues or mixing matrices.
\end{proposition}

\begin{conjecture}[Flavon hierarchy, P3$'$]\label{conj:flavon}
Sequential $SU(3)_F\to SU(2)_F\to U(1)_F$ breaking is driven by the
observer-indexed capacity RG flow $\omega^\star_\iota(\tau)$ on the flavon
potential $V_F(\phi_1,\phi_2,\phi_3)$, with $\phi_\iota$ the flavon along
diagonal $e_\iota$. The ordering
$\langle\phi_3\rangle\gg\langle\phi_2\rangle\gg\langle\phi_1\rangle$ follows
from the capacity weights $(1/27,10/27,16/27)$ of Theorem~\ref{thm:partition}
descended through the above--below correspondence
(Proposition~\ref{prop:above-below}), yielding Yukawa hierarchy and
CKM/PMNS mixing without an independent flavon postulate.
\end{conjecture}